\begin{longtable}{P{0.105\linewidth}P{0.39\linewidth}P{0.35\linewidth}}
\caption{Bounded observation, qualification and source anchor for each accepted study.}\\
\toprule ID & Extracted observation & Limitation and source location\\\midrule\endfirsthead
\toprule ID & Extracted observation & Limitation and source location\\\midrule\endhead
S001 \cite{reef} & GV100 PyTorch DISB-C: Orion p99 normalized latency1.17x vs Reef below1.05x; normalized total throughput0.97x vs1.14x(Fig29). TensorRT comparison trades throughput: mean1.21x Triton vs1.11x Reef(Fig28). All ratios retain model-solo normalization; do not mix engines/AMD strata. & Explicit OSDI2022 extension; select journal as preferred report and retain GA-S001 once. Requires deterministic/stateless kernels; closed-source comparison omits padding. Memory assumed sufficient; bandwidth contention and input-shape profiling limits. Printed publication date differs from some issue-year indexing; recorded2025 follows supplied copy and is in cutoff either way. \emph{Anchor:} Article pp1:23-1:24,1:31-1:34; sections6.2,6.6,7; Figs28-29; Markdown lines1111-1180,1600-1820 \\
S002 \cite{gpulet} & Local gpulets support interference-aware provisioning; cluster throughput is a different outcome. & Server-level gains cannot be treated as one-GPU utilization. \emph{Anchor:} PDF p10 Table2/3; section5; local profiling section3 \\
S003 \cite{pipeswitch} & Fig7: inference preemption latency reduced by up to 38.2ms. & Average switch/inference latency is not p99; multi-VM throughput separate. \emph{Anchor:} PDF pp11-13; sections6-6.2; Fig7 \\
S004 \cite{clockwork} & Fig5 measures goodput satisfying SLO, rather than raw request throughput. & Cross-hardware INFaaS comparison is indirect; controller/worker distinction. \emph{Anchor:} PDF p10; section6.1; Fig5 \\
S005 \cite{torpor} & BERT Table4: native42ms, remote43ms, PCIe swap144ms, NVLink swap45ms. & p98 tail; busy time is not SM occupancy; node scaling separate. \emph{Anchor:} PDF pp10-11; Table4; sections7.2-7.3 \\
S006 \cite{streambox} & Table2 evaluates runtime sandbox workflows on shared GPUs. & Sandbox memory separation is not proof of adversarial security. \emph{Anchor:} PDF p10 onward; section9; Table2 \\
S007 \cite{krypton} & Inference-specific sharing comparisons; section2.2 tests fault propagation. & Baseline adaptations and workload mix constrain attribution. \emph{Anchor:} PDF p10 Tables1-2; section5; section2.2 Fig3 \\
S008 \cite{sirius} & Inference p99/SLO outcomes accompany training progress; separate TTFT/TBT LLM case. & TaskSwitch is a reimplementation, not vanilla PipeSwitch. \emph{Anchor:} Section6.1 Fig10; section6.6 \\
S009 \cite{tally} & At 10\% idle time: normalized throughput Tally1.19 vs TGS1.65. & Latency protection sacrifices some BE throughput; author manuscript. \emph{Anchor:} PDF pp9-11; sections5-5.3; Fig6 \\
S010 \cite{parva} & Fig7 separates internal/external fragmentation; Fig5 counts GPUs. & Author manuscript; unusual runtime version; GPU count is not occupancy. \emph{Anchor:} PDF p8 onward; sectionIV; TableIV; Figs5/7 \\
S011 \cite{orion} & Interference-aware scheduling compares inference latency and BE throughput. & Prototype clients share a process; REEF-N differs from AMD REEF. \emph{Anchor:} PDF pp8-9; sections5.3-6.1; Tables1/3 \\
S012 \cite{gslice} & GUE measures efficiency; shadow instance loading takes about12s in Fig15. & GUE is not GPU busy percentage; median latency not p99; academic mirror. \emph{Anchor:} PDF pp8-12; section4; Table3; Fig15 \\
S013 \cite{krisp} & Kernel-wise right-sizing evaluates throughput and p95 inference latency. & Feature emulated; overhead subtracted; baseline is not NVIDIA MPS. \emph{Anchor:} PDF pp9-10; sectionsV-VI; TableIII \\
S014 \cite{xsched} & Fig15: relative p99 1.53/1.51/1.07 respectively versus solo. & Data assumed to fit memory; CUDA graphs and undocumented driver limits. \emph{Anchor:} PDF pp14-15; section8 Case3 Fig15; section9 \\
S015 \cite{paella} & p99 latency and throughput under bursty model mixes; Fig12 includes MPS. & Custom Triton TVM wrapper; MPS omitted from Fig11; dynamic batching unsupported. \emph{Anchor:} PDF pp9-12; section7 Table3/Figs11-12; section8 \\
S016 \cite{gnntriton} & Fig10: thrashing after two models/20GB or five/40GB. & Local slice evidence only; multi-server scaling separate; isolation properties described, not attack-tested. \emph{Anchor:} HTML Multi-model Scaling/Fig10; Triton implementation; limitations \\
S017 \cite{vastpipe} & A100 VGG16+GoogLeNet: throughput 2.7x no MPS and 1.5x default MPS; not a universal average or p99 result. NLP/CV mean 2.43x no MPS, 1.45x default MPS. MI100 mean 1.89x no mask, 1.24x default mask. & Include only local GPU inference evidence, not the main 1000-NPU cluster results. Offline Q-table assumes relatively stable request mix; GPU-specific quotas, repetitions and variance unresolved. Chart bars are not preserved in Markdown. Preliminary ICCD2021 report is linked as an unassessed version lead. \emph{Anchor:} Published pp457-459; sections6.7-6.9; Figs14-16; uploaded Markdown lines1150-1290 \\
S018 \cite{arkgpu} & Section5.2 reports all30 tested pairs satisfy the 2x-solo p99 QoS target with ArkGPU. Low-contention pairs can lose about5\% QPS from throttling. Four-GPU app-mix goodput gains94.98\% vs k8s-native and38.65\% vs MPS are separate system-level results. & Mixed inference/training qualifies via explicit inference p99/QPS and workload transfer rationale; it does not establish inference-only performance. Polling/API sleeps can add delay; BE checkpoint/restart affects costs. Huge goodput ratios can have near-zero baseline yield. Repetitions/variance and chart coordinates unresolved. \emph{Anchor:} Published pp313-319; sections4.3,5.1-5.4; Table3; Figs14-21; Markdown lines920-1360 \\
S019 \cite{dynap} & Single-A30 batch experiments: smaller partitions increase concurrency but can increase mean JCT; not online p99 evidence. Fig8 broader job set reports39.03\% lower makespan vs StaCon,12.14\% vs MISO and14.7\% higher normalized STP vs StaCon. Prediction-error sensitivity:15\% perturbation adds8.09\% makespan. & Retain one-GPU inference strata separately from two-GPU aggregate and cluster simulation. JCT/makespan/STP cannot replace request-tail latency or goodput. Reconfiguration can stop/restart active jobs; neither architectural MIG claims nor OOM avoidance prove hostile isolation. Model labels/parameter counts and corrupted overhead fraction require caution. \emph{Anchor:} PDF pp11-16; sections4.1-4.3.2; Table2; Figs8-13; Markdown lines700-865 \\
S020 \cite{igniter} & Local four-model latency prediction error1.53\%-5.02\%(Fig13). Fig14 deployment: iGniter6 GPUs/0 SLO violations vs gpu-lets+8/3,FFD+5/10,GSLICE+6/3;25\% lower hourly instance cost vs gpu-lets+, not25\% higher local utilization. Fig17 shows an initial p99 violation before shadow-process recovery. & Constant arrivals and four vision models limit bursty/LLM transfer. Negligible PCIe contention assumption can underpredict AlexNet; baselines modified. Steady-state/post-correction SLO success is not an unconditional guarantee; shadow standby incurs capacity/memory costs. Chart bars/latency traces unavailable in Markdown. \emph{Anchor:} Published pp820-824; sections4.2,5.1-5.4; Figs13-18; Markdown lines938-1245 \\
S021 \cite{inferfair} & Local Table3 load3:1:2: unfairness1.05 vs Clockwork1.50; SLO goodput524.29 vs457.24req/s. Other load ratios give smaller differences. Local unfairness is max/min model slowdown(Eq2), not p99 or security. Four-GPU MAF outcome kept separate. & Two-column Markdown interleaves algorithms/prose; use labelled Table3 rows and sections, not scrambled formulas. Repetitions/variance and software versions unresolved. MAF trace has minute-level invocation counts; cross-study batching and isolation definitions differ. \emph{Anchor:} Published pp17-18; sections5.1-5.4; Table3; Eq2; Markdown lines718-899 \\
S022 \cite{miger} & Section 3.1/Fig. 2 directly compares local MPS 57/29/14 percent and MIG 4g/2g/1g allocations. Section 5.2/Fig. 8 reports testbed mean JCT 46.6\% below MuxFlow and 36\% below MISO; these are job/cluster outcomes, not inference p99 gains. Fig. 9 notes slower small jobs under MIGER. & Includes mixed training/inference and multi-GPU optimization; retain only explicitly local sharing and inference-QoS evidence for assigned comparison. Abstract and introduction swap comparator order for the two JCT percentages; use section 5.2 mapping. No chart coordinates reconstructed or universal p99 claim. \emph{Anchor:} Printed pp. 505-506, 510-511; sections 3.1-3.3, 5.1-5.4; Tables 1-3; Figs. 2, 4, 8-11; N00001.md lines 119-229, 478-662 \\
S023 \cite{deft} & Sections 3.2 and 5 show optimal synchronization depends on batch and image size; the fixed strategy can impose much larger preempting-job JCT inflation after a workload change. Section 5.2 reports Deft switches to synchronizing every kernel for batch 30; earlier batch-1 case uses every 500 kernels. Table 2 cells are interleaved, so no exact violation-rate row is imported. & Preliminary six-page workshop evaluation with two principal models and one A40. Exhaustive pair profiling does not scale; no cluster/generative-model implementation claimed (section 6). Keep export year 2024 and printed conference/reference year 2023 separately; both inside window. \emph{Anchor:} Printed pp. 15-17; sections 3.2, 4.2, 5.1-5.3 and 6; Table 2 and Fig. 4 discussed in prose; N00005.md lines 250-468 \\
S024 \cite{fastgshare} & Local section 5.3 states eight RNNT pods at 12\% SM shares reach about 40 req/s with tail latency below 500 ms, versus one racing pod about 12 req/s and above 1250 ms. The local ResNet/RNNT comparison in Fig. 9 reports removing their observed interference with 24\% shares. No unlabelled tail value is relabelled p99. & Model-specific/proportional share effects, IPC memory trust and batching limit transfer. Prose fractional-gain notation is ambiguous (3.15x with raw 296.8/71.37), so no such multiplier imported. Supplied plot coordinates/percentile identity and repetitions not reconstructed. \emph{Anchor:} Printed pp. 641-643; sections 4, 5.1-5.4; Figs. 8-10; N00015.md lines 523-735 \\
S025 \cite{split} & Section 5.5 reports violation rate below 10\% beyond alpha=4 versus 26\% under RT-A in the discussed comparison. For short requests, jitter reductions versus Clockwork/PREMA/RT-A are 55.3/46.8/68.9\% at low load and 56.0/50.3/69.3\% at high load. Jitter is block-execution standard deviation, not p99 latency. Section 6 reports split overhead and static-graph limits. & One edge GPU and synthetic Poisson load; not a server-class GPU generalization. Temporal predictability sacrifices spatial parallelism; long-request stability can worsen. Table 3 split overhead and RNN limits qualify short-request improvements. \emph{Anchor:} Printed pp. 611-613; sections 5.1-5.5, 6; Tables 1-3; Figs. 6-7; N00040.md lines 451-629 \\
S026 \cite{migserve} & Section 5.2 Azure: sustainable rate at 90\% SLO attainment is 9.84 RPS versus LA-Dyn 4.13 and Static 1.85. Arena approximately 10.99 versus 2.37 and 2.48 RPS. These are four-GPU server outcomes; local mechanism is MIG layout control, not a single-GPU throughput multiplier. & Four-GPU evaluation and constructed layout-agnostic comparator; 90\% attainment is weaker than 99\% SLO targets used elsewhere. MIG hardware isolation is described, without adversarial/fault experiment. Source pixels and unreported repeats/versions are not reconstructed. \emph{Anchor:} Printed pp. 850-853; sections 4, 5.1-5.4; Fig. 9 and Figs. 11-13; N00831.md lines 519-780 \\
S027 \cite{mmk} & Section 4.3/Fig. 9 reports mean JCT 28\% lower, makespan 32\% lower and normalized progress-rate STP 35\% higher than MIGER. These are mixed-job/testbed metrics, not per-request p99 latency. Section 4.5 reports 1.5\% MIG and 3.2\% MPS reconfiguration lifecycle overhead. & Mixed inference/training and two-GPU optimization; local priority scheduling is the transferable mechanism. Many figure labels contain NUL/glyph corruption, so numerical prose only. STP/JCT does not establish online tail-latency or fault containment. \emph{Anchor:} Article pp. 84:15-84:19; sections 4.1-4.5; Table 3; Figs. 9-12; N00839.md lines 789-1040 \\
S028 \cite{boer} & Local section 2.4/Fig. 3 observes ParvaGPU can underperform MIG-only serving for larger BERT/ResNet50 while helping lightweight MobileNet. Section 5.2/Fig. 8 gives eight-GPU QoS-preserved throughput mean 34.61\% above ParvaGPU; 56.8/62.49/34.63/44.83\% above MIG-serving/gpulet/USHER/PARIS. Keep local qualitative counterexample separate from fleet improvements. & Bayesian exploration cost and instance interference qualify hybrid gains. Baseline reductions (MIG-serving fast/USHER IR only) and multi-GPU request fan-out constrain comparison. Same senior author as MIGER, but different inference-only system and experiments; not a same-study report extension identified in supplied text. \emph{Anchor:} Printed pp. 521, 526-528; sections 2.4, 4.6, 5.1-5.4; Tables 2-3; Figs. 3, 8, 10; N00844.md lines 200-265, 700-890 \\
S029 \cite{admix} & Section 4.2 reports average best-effort throughput 1.58x Static MPS with local LLaMA-7B (TP1) and 3.9x with Qwen-14B (TP2), keeping strata distinct. AdMix still increases mean TTFT/TPOT versus LLM-only; on Qwen its mean latency increases exceed Static MPS. More throughput is not zero interference. & Full six-page original workshop report; one/two-GPU workload strata and tuned-on-test-trace static reference limit transfer. Estimates of average latency do not establish p99. No testbed code rerun or security/fault experiment inferred. \emph{Anchor:} Printed pp. 3-5; sections 3 and 4.1-4.2; Table 1; Figs. 5-6; N00872.md lines 270-453 \\
S030 \cite{sgdrc} & Section 9.3: versus Orion, P40 overall throughput up to 1.47x and BE up to 2.36x; A2000 up to 1.32x and 1.79x, respectively. Mean SLO attainment 99.0\% across evaluated settings. Section 10 explicitly states shared-context GPU runtime errors are not isolated. & Software coloring is not a proven security boundary. MIG integration is future work; not benchmarked on A100/H100. Reverse-engineering/calibration and TVM compatibility matter; do not treat throughput maxima as averages or infer a P40 MPS result. \emph{Anchor:} Printed pp. 276-277; sections 9.2-9.3, 10; Table 3 and Fig. 18; N00904.md lines 863-965 \\
S031 \cite{enertune} & Section 2.1/Fig. 2: co-location choice can increase energy; SDXL plus MobileNetV3 example uses 13\% less energy versus no sharing despite 1.14x latency. Section 7.1/Fig. 9: 16-GPU testbed energy reduced by factors 1.4-2.3 and power by 1.3-2.6 versus baselines while meeting SLOs; fleet results are not local per-GPU savings. & Joint power/energy objectives select different allocations. Local model-pair evidence must stay distinct from fleet savings and from frequency-only effects. Source chart pixels not available; experiments/code not rerun; missing warm-up and independent-run uncertainty not assumed. \emph{Anchor:} Printed pp. 66-67, 73-75; sections 2.1-2.2, 6, 7.1-7.2; Figs. 2, 9; N00946.md lines 140-249, 798-995 \\
S032 \cite{fluidfaas} & Section 7.1/Fig. 9 reports FluidFaaS SLO hit-rate improvement of 90\% for medium and 61\% for heavy workload mixes over the compared systems; the abstract reports 25--75\% throughput increases across workloads. These are system-level results, not a measured per-GPU utilization effect. & Pipeline, placement and MIG allocation contribute jointly; comparisons with a non-pipelined baseline cannot isolate the benefit of MIG partitioning alone. No adversarial security or fault-injection evidence is inferred from MIG descriptions. \emph{Anchor:} N00920.md abstract, sections 3, 6, 7.1; Tables 2, 5--6; Figs. 9 and 12 \\
S033 \cite{bless} & The introduction and evaluation report mean latency reductions of 39.1\% for two BERT inference instances and 41.2\% for four, under the paper's compared GPU-sharing systems. Other reported aggregate gains are not recast as p99 or universal speedups. & Quota, model pairing and GPU generation condition results. The conversion has damaged column/table layout, so no unlabeled plotted coordinates or precise tail-percentile gains were recovered. \emph{Anchor:} N01085.md intro, sections 4, 6.1--6; Fig. 1 and BERT co-location result prose \\
S034 \cite{interference_deeper} & H100 Llama decode P90 time between tokens rises from 7.53 to 16.56 ms under a sleep co-tenant in section 4.2 (batch 1, 1,000-token prompt). Table 1's Green Context experiment isolates SM sets (64/68) but reports P90 TBT as high as 18.38 ms under copy-kernel memory pressure; its alone baseline is 14.2 ms on the whole GPU, so partition cost and interference cannot be collapsed into one effect. & Synthetic co-tenants and short decode traces quantify resource contention; the paper proposes a future scheduler but does not implement an end-to-end serving controller. Green Context SM partition does not partition L2 or external memory bandwidth. Table 1 conversion is interleaved; only prose and clearly attributable values are used. \emph{Anchor:} N01096.md sections 3--5, especially 4.1--4.4; Fig. 2 and Table 1 \\
S035 \cite{edge_sky} & Abstract and multi-tenancy evaluation report inference throughput gains under concurrent execution; do not attribute heterogeneous placement gains solely to GPU sharing. & Edge GPU, CPU and TPU contributions are coupled; throughput is not latency predictability or security isolation. \emph{Anchor:} N00027.md abstract; multi-tenancy and performance evaluation sections, figures on CME and DMP \\
S036 \cite{miriam} & Abstract reports up to 92\% throughput improvement while critical tasks incur less than 10\% latency overhead in its evaluated scenarios. & Edge elastic-kernel scheduling is transferable temporal/spatial evidence, not a direct MIG/MPS comparison. \emph{Anchor:} N00046.md abstract; sections on kernel coordination and evaluation \\
S037 \cite{elasticroom} & The evaluated engine balances request goodput, throughput and latency under co-located inference; comparison is engine-level. & Compiler and priority scheduling effects cannot be isolated as pure partition effects; plotted numbers were not transcribed. \emph{Anchor:} N00048.md abstract; section 6.2 evaluation and Fig. 1 \\
S038 \cite{switchflow} & Section 5.4/Fig. 10 reports interleaved execution gains among inference jobs of up to 30\% versus session-based time slicing in tested pairings. & The 19.05x p95 inference latency example in Fig. 6 uses a training co-tenant and is not evidence for two inference tenants. \emph{Anchor:} N00049.md sections 5.1, 5.2.1 and 5.4; Figs. 6 and 10 \\
S039 \cite{itif} & Abstract reports up to 2.40x performance versus baselines for evaluated shared-backbone transformer inference. & Shared backbone changes model computation; speedup cannot be assigned to GPU scheduling alone or generalized to unrelated models. \emph{Anchor:} N00059.md abstract; design and experiment sections including four-tenant setup \\
S040 \cite{protean} & The evaluation reports SLO compliance and cost effects of orchestration under trace-driven inference; no per-GPU improvement is extracted from system-level totals. & Multiple GPUs and global scheduling contribute jointly; local mechanism claims cannot be isolated without matched per-device comparisons. \emph{Anchor:} N00063.md abstract, sections on MIG/MPS design and evaluation of 22 workloads \\
S041 \cite{wattwiser} & Section 4/Fig. 3 shows model tail latency relative to SLO while requests consolidate onto shared GPUs; power savings are whole-server effects. & The four-GPU power outcome cannot estimate per-GPU partition efficiency; MI50 setup is not a MIG evaluation. \emph{Anchor:} N00068.md sections 4.1--4.2 and Fig. 3 \\
S042 \cite{deinfer} & Section 5 evaluates prediction accuracy and allocation for concurrent inference with MPS; its Fig. 2 shows that MPS spatial limits do not eliminate shared-resource latency interference. & Two GPUs are present in testbed; claims here concern within-device workloads, not a two-device pooled speedup. \emph{Anchor:} N00070.md sections 2.2 and 5.1--5.2; Fig. 2, Table 1 and testbed table \\
S043 \cite{esg} & Abstract reports 61--80\% SLO hit-rate improvement and 47--187\% cost savings over prior systems for evaluated workflows. & Figures are workflow/system results influenced by batching and scheduler search, not isolated effects of local GPU sharing. \emph{Anchor:} N00073.md abstract, sections 3 and 5; evaluation setup and end-to-end comparison \\
\bottomrule\end{longtable}
